\chapter{Architetture software moderne}
% Il codice in linea (\code) non va a capo: un po' di elasticita' in piu' evita righe che sbordano
\emergencystretch=2.5em

Il capitolo precedente ha presentato gli stili architetturali classici: client/server, repository, livelli, three-tier,
SOA. Qui guardiamo le architetture dei sistemi \textbf{moderni}, che devono servire milioni di utenti, non fermarsi mai ed
evolvere ogni giorno. Il filo conduttore sono le \textbf{proprietà} che l'architettura determina (scalabilità,
disponibilità, consistenza dei dati) e le tecniche per ottenerle: bilanciamento del carico, replica, code di messaggi,
microservizi. Chiudiamo con un tema meno tecnico ma altrettanto decisivo: il legame fra l'architettura e
l'\textbf{organizzazione} dei team.

% Stili delle figure di questo capitolo
\tikzset{
  pc/.style={rectangle, rounded corners=2pt, draw=swteal!80!black, fill=swteal!10, line width=0.7pt,
             minimum width=#1, minimum height=#1, font=\scriptsize, inner sep=0pt},
  pc/.default=0.7cm,
  servizio/.style={rectangle, rounded corners=2pt, draw=sworange!80!black, fill=sworange!12, line width=0.8pt,
             minimum width=1.5cm, minimum height=0.8cm, font=\scriptsize\bfseries, align=center},
  db/.style={cylinder, shape border rotate=90, aspect=0.22, draw=swblue!70!black, fill=swblue!14, line width=0.8pt,
             minimum width=1.2cm, minimum height=0.9cm, font=\scriptsize\bfseries, align=center},
  fr/.style={-{Stealth[length=2.2mm]}, line width=0.8pt, draw=swgray},
  frd/.style={{Stealth[length=2.2mm]}-{Stealth[length=2.2mm]}, line width=0.8pt, draw=swgray},
}

% ============================================================
\section{Che cos'è, di nuovo, un'architettura}

Nel Capitolo 11 abbiamo definito l'architettura come l'insieme delle strutture di un sistema. Martin Fowler ne dà una
definizione complementare, centrata sulle \textbf{decisioni}.

\dfn{Architettura software (Fowler)}{
  L'architettura di un sistema software è l'insieme delle \textbf{decisioni di progetto principali} sul sistema: il
  progetto di massima (\textit{blueprint}) della sua costruzione e della sua evoluzione. Sono decisioni che hanno un
  impatto significativo sulla struttura del sistema, sul suo comportamento e sulle sue interazioni, e sulle sue proprietà
  non funzionali.
}

«Principale» indica un grado di importanza tale da dare a una decisione lo \textbf{status di decisione architetturale}.
Ma quale lo è?
\begin{itemize}
  \item Scegliere fra MariaDB e MySQL è una decisione architetturale? Probabilmente no: sono due database relazionali quasi
        intercambiabili, e passare dall'uno all'altro costa poco.
  \item Scegliere fra un database \textbf{relazionale} e uno \textbf{NoSQL}? Sì: cambia il modo in cui si modellano i dati,
        le garanzie di consistenza, il codice che li usa. Tornare indietro costa moltissimo.
\end{itemize}

L'architettura riguarda ciò che è \textbf{difficile da cambiare in seguito}. Non esiste un criterio oggettivo per stabilire
quali decisioni siano «principali»; come scrive Fowler, \textit{«l'architettura riguarda le cose importanti, qualunque esse
siano»}. Il cuore del pensare un software in modo architetturale è decidere che cosa è importante, e poi spendere energie
per mantenere quegli elementi in buone condizioni. Per diventare architetto, uno sviluppatore deve saper riconoscere gli
elementi importanti: quelli che, se non tenuti sotto controllo, porteranno con ogni probabilità a problemi seri.

\subsection{Perché l'architettura conta}

Per clienti e utenti l'architettura è un argomento ostico: è qualcosa che \textbf{non percepiscono} direttamente. Eppure ha
un impatto significativo sulla struttura del sistema, sul suo comportamento e sulle sue proprietà non funzionali. Inoltre
un'architettura scadente rende il sistema molto più difficile da sviluppare e da far evolvere, con conseguenze sui costi e
sui tempi di consegna delle nuove funzionalità.

\Needspace{18\baselineskip}
\noindent\begin{minipage}[c]{0.44\textwidth}
Il grafico, da un articolo di Fowler (\textit{Is High Quality Software Worth the Cost?}), mostra la funzionalità
accumulata nel tempo da due progetti. Quello con \textbf{bassa qualità interna} parte più veloce, ma rallenta presto:
ogni nuova funzionalità costa di più, perché il codice è sempre più difficile da capire e modificare. Quello con
\textbf{alta qualità interna} (una buona architettura) parte un po' più lento, ma poi consegna più in fretta e a costi
minori.

Il punto chiave è \textbf{quando} le due curve si incrociano: secondo Fowler, dopo \textbf{settimane}, non mesi. La qualità
interna costa di più all'inizio, ma ripaga molto presto.
\end{minipage}\hfill
\begin{minipage}[c]{0.53\textwidth}
\centering
\begin{tikzpicture}[x=0.75cm, y=0.75cm]
  \draw[-{Stealth[length=2mm]}, line width=0.8pt] (0,0) -- (10.2,0) node[below left, note]{tempo};
  \draw[-{Stealth[length=2mm]}, line width=0.8pt] (0,0) -- (0,6.4) node[above right, note, align=left]{funzionalità\\ accumulata};
  \draw[line width=1.3pt, draw=swblue, domain=0:9.8, samples=60] plot (\x, {5.6*(1-exp(-\x/5.5))*(1+0.04*\x)});
  \draw[line width=1.3pt, draw=q9!70, domain=0:9.8, samples=60] plot (\x, {3.2*(1-exp(-\x/1.8))});
  \draw[dashed, draw=swgray] (2.95,0) -- (2.95,2.6);
  \node[note, anchor=north, align=center, text=q4!80!black] at (2.95,-0.05) {dopo settimane,\\ non mesi};
  \node[note, anchor=west, text=swblue, font=\scriptsize\bfseries] at (3.9,5.9) {alta qualità interna};
  \node[note, anchor=west, text=q9!80!black, font=\scriptsize\bfseries] at (6.9,3.55) {bassa qualità interna};
  \node[note, anchor=west, align=left, text=swgray, font=\tiny] at (0.4,3.6) {all'inizio l'alta qualità\\ rallenta un poco};
\end{tikzpicture}
\end{minipage}

\subsection{Le proprietà che dipendono dall'architettura}

Le proprietà su cui l'architettura incide di più sono sei. Le prime tre sono i temi delle prossime sezioni; le altre tre
sono misure di prestazione che conviene distinguere bene.

\begin{center}
\begin{tikzpicture}[prop/.style={chip=#1, minimum width=3.0cm, minimum height=0.7cm, font=\small\bfseries}]
  \node[prop=swteal] at (0,0) {Scalabilità};
  \node[prop=swteal] at (3.4,0) {Disponibilità};
  \node[prop=swteal] at (6.8,0) {Consistenza};
  \node[prop=swblue] at (0,-0.95) {Prestazioni};
  \node[prop=swblue] at (3.4,-0.95) {Latenza};
  \node[prop=swblue] at (6.8,-0.95) {Throughput};
\end{tikzpicture}
\end{center}

\begin{itemize}
  \item La \term{latenza} è il \textbf{tempo} necessario a servire una singola richiesta (ad esempio 200 ms).
  \item Il \term{throughput} è la \textbf{quantità} di lavoro servita nell'unità di tempo (ad esempio 5000 richieste al
        secondo).
\end{itemize}
Le due misure sono legate ma non coincidono: un sistema può servire moltissime richieste al secondo, ciascuna lentamente.

% ============================================================
\section{Scalabilità}

\dfn{Scalabilità}{
  Un servizio è \term{scalabile} se, aggiungendo risorse, le sue prestazioni aumentano in modo \textbf{proporzionale} alle
  risorse aggiunte.
}

Di solito aumentare le prestazioni significa servire \textbf{più unità di lavoro} (più richieste, più utenti), ma può anche
significare gestire unità di lavoro \textbf{più grandi}, ad esempio quando i dati crescono. Un modo semplice di distinguere
prestazioni e scalabilità:
\begin{itemize}
  \item se c'è un problema di \textbf{prestazioni}, il sistema è lento anche per un \textbf{singolo} utente;
  \item se c'è un problema di \textbf{scalabilità}, il sistema è veloce per un singolo utente, ma lento \textbf{sotto carico}.
\end{itemize}

\Needspace{16\baselineskip}
\noindent\begin{minipage}[c]{0.5\textwidth}
\textbf{La scalabilità è una questione di soldi.} Il grafico mostra come crescono i costi al crescere del volume di lavoro
(richieste servite, transazioni, utenti contemporanei\ldots) per tre ipotetici sistemi.
\begin{itemize}
  \item \textcolor{swred}{\textbf{Male}}: costi che crescono più che proporzionalmente. Il sistema è economico finché è
        piccolo, ma diventa insostenibile proprio quando ha successo.
  \item \textcolor{swteal}{\textbf{Accettabile}}: costi che crescono in proporzione al volume.
  \item \textcolor{swgreen}{\textbf{Bene}}: costi che crescono meno che proporzionalmente. Ogni utente in più costa meno
        del precedente.
\end{itemize}
Le tre curve si incrociano in un punto: a quel volume i tre sistemi costano uguale. Prima di quel punto il sistema
«cattivo» è addirittura il più economico, ed è proprio questa la trappola: le differenze si vedono solo quando il volume
cresce.
\end{minipage}\hfill
\begin{minipage}[c]{0.46\textwidth}
\centering
\begin{tikzpicture}[x=0.7cm, y=0.7cm]
  \draw[-{Stealth[length=2mm]}, line width=0.8pt] (0,0) -- (8.6,0) node[below left, note]{volume};
  \draw[-{Stealth[length=2mm]}, line width=0.8pt] (0,0) -- (0,6.6) node[above right, note]{costi};
  \draw[line width=1.3pt, draw=swred, domain=0:4.6, samples=40] plot (\x, {0.8+0.2444*\x*\x});
  \draw[line width=1.3pt, draw=swteal, domain=0:8.2, samples=10] plot (\x, {1.2+0.6*\x});
  \draw[line width=1.3pt, draw=swgreen, domain=0:8.2, samples=40] plot (\x, {1.6+1.01*ln(1+\x)});
  \node[note, text=swred, font=\scriptsize\bfseries, anchor=south] at (4.6,6.3) {male};
  \node[note, text=swteal, font=\scriptsize\bfseries, anchor=west] at (8.2,6.3) {accettabile};
  \node[note, text=swgreen, font=\scriptsize\bfseries, anchor=west] at (8.2,3.85) {bene};
  \draw[dashed, draw=swgray] (3,0) -- (3,3);
  \fill[black] (3,3) circle (2pt);
  \node[note, anchor=north, font=\tiny] at (3,-0.05) {stesso costo};
\end{tikzpicture}
\end{minipage}

\subsection{Scalare orizzontalmente, verticalmente, diagonalmente}

Come si serve più lavoro? Le strade sono tre.

\begin{figure}[H]
\centering
\begin{tikzpicture}
  \tikzset{riq/.style={rectangle, rounded corners=4pt, draw=#1!70!black, fill=#1!8, line width=0.8pt,
                       minimum width=4.6cm, minimum height=4.2cm}}
  % orizzontale
  \node[riq=q5] at (0,0) {};
  \foreach \r in {0,1,2} \foreach \c in {0,1,2} \node[pc] at (-1.0+\c*1.0,1.3-\r*0.95) {\icpu};
  \draw[fr, draw=q5!80!black, line width=1.2pt] (-1.6,-1.45) -- (1.6,-1.45);
  \node[font=\scriptsize\bfseries, align=center] at (0,-2.55) {Scalabilità orizzontale};
  \node[note, align=center] at (0,-3.0) {più nodi di calcolo};
  % verticale
  \node[riq=q6] at (5.0,0) {};
  \node[pc=0.55cm] at (5.0,-1.35) {\icpu};
  \node[pc=0.9cm] at (5.0,-0.4) {\icpu};
  \node[pc=1.3cm, font=\large] at (5.0,1.05) {\icpu};
  \draw[fr, draw=q6!80!black, line width=1.2pt] (3.25,-1.6) -- (3.25,1.7);
  \node[font=\scriptsize\bfseries, align=center] at (5.0,-2.55) {Scalabilità verticale};
  \node[note, align=center] at (5.0,-3.0) {un nodo più potente};
  % diagonale
  \node[riq=q7] at (10.0,0) {};
  \node[pc=0.5cm] at (8.6,-1.4) {\icpu};
  \node[pc=0.65cm] at (8.6,-0.5) {\icpu}; \node[pc=0.65cm] at (9.4,-0.5) {\icpu};
  \node[pc=0.85cm] at (8.7,0.6) {\icpu}; \node[pc=0.85cm] at (9.7,0.6) {\icpu}; \node[pc=0.85cm] at (10.7,0.6) {\icpu};
  \draw[fr, draw=q7!80!black, line width=1.2pt] (9.5,-1.5) -- (11.7,-0.2);
  \node[font=\scriptsize\bfseries, align=center] at (10.0,-2.55) {Scalabilità diagonale};
  \node[note, align=center] at (10.0,-3.0) {più nodi, e più potenti};
\end{tikzpicture}
\caption{Le tre strade per servire più lavoro.}
\end{figure}

La scalabilità \textbf{verticale} è la più semplice (non cambia nulla nel software), ma ha un limite fisico ed economico:
oltre una certa soglia una macchina più potente costa moltissimo, o non esiste. Quella \textbf{orizzontale} non ha limiti
teorici, ma richiede che il software sappia distribuire il lavoro su più nodi.

\subsection{Il bilanciamento del carico}

Quando si scala orizzontalmente (o diagonalmente), serve un componente dedicato che distribuisca le richieste in arrivo fra
i nodi di calcolo: il \term{load balancer} (bilanciatore del carico). Le strategie più comuni per scegliere il nodo sono:
\begin{itemize}
  \item \textbf{casuale}: ogni richiesta va a un nodo scelto a caso;
  \item \textbf{meno carico} (\textit{least loaded}): la richiesta va al nodo meno occupato in quel momento;
  \item \textbf{round robin}: i nodi si alternano a turno, uno dopo l'altro.
\end{itemize}

\Needspace{10\baselineskip}
\noindent\begin{minipage}[c]{0.46\textwidth}
\centering
\begin{tikzpicture}
  \node[note, font=\scriptsize] (w) at (0,0) {richieste};
  \node[rectangle, rounded corners=3pt, draw=swblue!70!black, fill=swblue!15, line width=0.8pt, font=\scriptsize\bfseries,
        align=center, minimum height=0.9cm] (lb) at (2.4,0) {load\\ balancer};
  \foreach \i/\y in {1/1.1, 2/0, 3/-1.1} {
    \node[pc] (n\i) at (5.0,\y) {\icpu};
    \draw[fr] (lb.east) -- (n\i.west);
  }
  \draw[fr] (w) -- (lb);
\end{tikzpicture}
\end{minipage}\hfill
\begin{minipage}[c]{0.5\textwidth}
\warning{{Il load balancer è un punto delicato}}{
  Tutto il traffico passa da lì, quindi può diventare un \textbf{collo di bottiglia}. E un singolo load balancer è un
  \textbf{single point of failure}: se si guasta, si ferma tutto, anche se i nodi dietro funzionano. Nei sistemi reali i
  load balancer stessi vengono replicati.
}
\end{minipage}

% ============================================================
\section{Disponibilità}

\dfn{Disponibilità}{
  La \term{disponibilità} (\textit{availability}) è la \textbf{frazione di tempo} in cui un sistema o un servizio è
  operativo e accessibile.
}

La disponibilità è cruciale per i sistemi che devono funzionare in modo continuo: i sistemi \textbf{business-critical},
dove un fermo causa perdite economiche e danni d'immagine, e i sistemi \textbf{safety-critical}, dove sono in gioco persone.
Nel modello ISO/IEC 25010 del Capitolo 4 è una delle sotto-caratteristiche dell'affidabilità. Un sistema può diventare
indisponibile per:
\begin{itemize}
  \item problemi \textbf{software} (l'applicazione smette di rispondere);
  \item \textbf{carichi eccessivi} (non ci sono abbastanza risorse di calcolo);
  \item perdita della \textbf{connessione di rete};
  \item \textbf{guasti hardware}.
\end{itemize}

\subsection{Quantificare la disponibilità: i «nove»}

La disponibilità si esprime come percentuale di tempo in cui il servizio è attivo (\textit{uptime}), e si misura in
\textbf{numero di nove}: un servizio disponibile al 99{,}99\% ha «quattro nove». Ogni nove in più divide per dieci il fermo
annuo tollerato.

\begin{center}
\renewcommand{\arraystretch}{1.3}
\begin{tabular}{@{}llr@{}}
  \toprule
  \textbf{Nove} & \textbf{Disponibilità} & \textbf{Fermo tollerato in un anno} \\
  \midrule
  uno     & 90\%        & circa 36 giorni e mezzo \\
  due     & 99\%        & circa 3 giorni e 15 ore \\
  tre     & 99{,}9\%    & circa 8 ore e 45 minuti \\
  quattro & 99{,}99\%   & circa 52 minuti \\
  cinque  & 99{,}999\%  & circa 5 minuti \\
  sei     & 99{,}9999\% & circa 31 secondi \\
  \bottomrule
\end{tabular}
\end{center}

Il calcolo è semplice: un anno ha circa $365 \times 24 = 8760$ ore, e il fermo tollerato è la frazione restante. Con tre
nove, $8760 \times 0{,}001 = 8{,}76$ ore, cioè circa 8 ore e 45 minuti (le slide riportano 8 ore e 41 minuti: la piccola
differenza dipende da come si conta la durata dell'anno).

\subsection{Come si ottiene un'alta disponibilità: il fail-over}

L'idea è non dipendere da un'unica istanza di un componente. Le tecniche di \term{fail-over} sono due.

\Needspace{14\baselineskip}
\noindent\begin{minipage}[t]{0.49\textwidth}
\gbox{Fail-over attivo-passivo}{swblue}{\raggedright
Se il componente \textbf{primario} si guasta, un'istanza \textbf{secondaria} ne prende il posto. La secondaria di solito è
\textbf{passiva}: non fa nulla finché il primario funziona.

\centering
\begin{tikzpicture}
  \node[note] (r) at (0,0) {richieste};
  \node[pc] (p) at (2.6,0.55) {\icpu};
  \node[pc, draw=swgray!60, fill=swlight, text=swgray!60] (s) at (2.6,-0.55) {\icpu};
  \draw[fr] (r.east) -- (p.west);
  \draw[fr, dashed, draw=swgray!60] (r.east) -- (s.west);
  \node[note, anchor=west] at (3.0,0.55) {primario (lavora)};
  \node[note, anchor=west, text=swgray] at (3.0,-0.55) {secondario (in attesa)};
\end{tikzpicture}}
\end{minipage}\hfill
\begin{minipage}[t]{0.49\textwidth}
\gbox{Fail-over attivo-attivo}{swgreen}{\raggedright
Esistono fin dall'inizio due o più istanze, che si dividono il lavoro (come nella scalabilità orizzontale). Se una si
guasta, le altre si fanno carico anche della sua parte.

\centering
\begin{tikzpicture}
  \node[note] (r) at (0,0) {richieste};
  \node[pc] (p) at (2.6,0.55) {\icpu};
  \node[pc] (s) at (2.6,-0.55) {\icpu};
  \draw[fr] (r.east) -- (p.west);
  \draw[fr] (r.east) -- (s.west);
  \node[note, anchor=west] at (3.0,0.55) {lavora};
  \node[note, anchor=west] at (3.0,-0.55) {lavora anche lei};
\end{tikzpicture}}
\end{minipage}

\subsection{Osservabilità e monitoraggio}

Con il fail-over, la durata del fermo dipende da quanto in fretta si reagisce al guasto.
\begin{itemize}
  \item \textbf{Quanto in fretta ci accorgiamo che un componente si è guastato?} E, prima ancora, \textbf{come} ce ne
        accorgiamo? Di solito con gli \term{heartbeat}: segnali o controlli periodici che verificano che il sistema funzioni
        normalmente. Se un componente smette di «battere», lo si considera guasto.
  \item \textbf{Quanto ci mette la secondaria a partire?} Può essere in \textit{hot standby} (già avviata e pronta a
        subentrare) o in \textit{cold standby} (spenta, da avviare al momento). Nel fail-over attivo-attivo il problema non
        si pone: le altre istanze sono già in funzione.
\end{itemize}

\dfn{Osservabilità}{
  L'\term{osservabilità} è il grado in cui riusciamo a capire che cosa sta succedendo all'interno di un sistema software
  mentre è in esecuzione, a partire dai dati che produce (log, metriche, tracce delle richieste). È una caratteristica
  importante di un'architettura; strumenti come \textit{OpenTelemetry} servono proprio a raccogliere questi dati in modo
  uniforme.
}

% ============================================================
\section{Dati e consistenza}

Scalare orizzontalmente le istanze dell'applicazione è stato relativamente facile: basta un nuovo server con
l'applicazione e un load balancer davanti. Ma quasi tutte le applicazioni hanno anche dati \textbf{persistenti}, di solito in
un database relazionale. Tutte le istanze, dietro il load balancer, usano lo stesso database: e se è il database a non
reggere più il carico? Scalare i dati è molto meno semplice.

\begin{center}
\begin{tikzpicture}
  \node[note] (w) at (0,0) {richieste};
  \node[rectangle, rounded corners=3pt, draw=swblue!70!black, fill=swblue!15, line width=0.8pt, font=\scriptsize\bfseries,
        align=center, minimum height=0.9cm] (lb) at (2.4,0) {load\\ balancer};
  \foreach \i/\y in {1/1.1, 2/0, 3/-1.1} {
    \node[pc] (n\i) at (5.0,\y) {\icpu};
    \draw[fr] (lb.east) -- (n\i.west);
    \draw[fr] (n\i.east) -- (7.3,0);
  }
  \draw[fr] (w) -- (lb);
  \node[db] (d) at (8.0,0) {DB};
  \node[note, anchor=west, text=swred, align=left] at (8.8,0) {un solo database\\ per tutte le istanze?};
\end{tikzpicture}
\end{center}

\subsection{Le proprietà ACID}

\dfn{Proprietà ACID}{
  Le transazioni dei database relazionali (RDBMS) godono di quattro proprietà.
  \begin{itemize}
    \item \textbf{Atomicità} (\textit{Atomicity}): ogni transazione è «tutto o niente».
    \item \textbf{Consistenza} (\textit{Consistency}): ogni transazione porta il database da uno stato valido a un altro
          stato valido.
    \item \textbf{Isolamento} (\textit{Isolation}): eseguire transazioni in modo concorrente dà lo stesso risultato che
          eseguirle in sequenza.
    \item \textbf{Durabilità} (\textit{Durability}): una volta confermata (\textit{commit}), una transazione resta tale,
          anche in caso di guasto.
  \end{itemize}
}

Sono proprio queste garanzie a rendere impossibile scalare un RDBMS orizzontalmente in modo semplice: se i dati stanno su
più macchine, garantire isolamento e consistenza richiede che le macchine si coordinino a ogni scrittura. Le strade sono:
\begin{itemize}
  \item scalarlo \textbf{verticalmente}, dandogli più risorse di calcolo;
  \item usare tecniche specifiche: \textbf{replica} (master-slave e master-master), \textbf{federazione} e
        \textbf{sharding}, \textbf{denormalizzazione}.
\end{itemize}

\subsection{Replica master-slave e master-master}

\begin{figure}[H]
\centering
\begin{tikzpicture}[
    cl/.style={rectangle, rounded corners=2pt, draw=black!70, fill=white, line width=0.6pt, font=\scriptsize,
               minimum width=1.2cm, minimum height=0.55cm},
    ms/.style={rectangle, draw=black!70, fill=#1, line width=0.7pt, font=\scriptsize\bfseries, minimum width=1.3cm,
               minimum height=0.6cm},
    ms/.default=black!12,
    lab/.style={font=\tiny, above, inner sep=1pt}]
  % master-slave
  \node[font=\small\bfseries] at (2.6,1.35) {Master-slave};
  \node[cl] (c1) at (0,0.6) {client};
  \node[cl] (c2) at (0,-0.6) {client};
  \node[cl] (c3) at (0,-1.6) {client};
  \node[ms] (m) at (3.6,0.6) {master};
  \node[ms=white] (s1) at (2.6,-0.6) {slave};
  \node[ms=white] (s2) at (4.6,-1.6) {slave};
  \draw[dashed, draw=swgray] (1.6,1.05) rectangle (5.6,-2.1);
  \draw[fr] (c1) -- node[lab]{letture/scritture} (m);
  \draw[fr] (c2) -- node[lab]{letture} (s1);
  \draw[fr] (c3) -- node[lab]{letture} (s2);
  \draw[fr, draw=swred] (m) -- node[font=\tiny, left]{replica} (s1);
  \draw[fr, draw=swred] (m) -- (s2);
  % master-master
  \begin{scope}[shift={(8.2,0)}]
    \node[font=\small\bfseries] at (2.2,1.35) {Master-master};
    \node[cl] (d1) at (0,0.4) {client};
    \node[cl] (d2) at (0,-1.3) {client};
    \node[ms] (m1) at (3.4,0.4) {master};
    \node[ms] (m2) at (3.4,-1.3) {master};
    \draw[dashed, draw=swgray] (1.6,1.05) rectangle (5.2,-2.1);
    \draw[fr] (d1) -- node[lab]{letture/scritture} (m1);
    \draw[fr] (d2) -- node[lab]{letture/scritture} (m2);
    \draw[frd, draw=swred] (m1) -- node[font=\tiny, right]{sincronizzazione} (m2);
  \end{scope}
\end{tikzpicture}
\caption{Le due forme di replica (rielaborate dalle slide \textit{Scalability, availability, stability, patterns}). In
         rosso il lavoro di coordinamento che introduce un \textbf{sovraccarico}.}
\end{figure}

\begin{itemize}
  \item \textbf{Master-slave.} Ci sono più istanze del RDBMS: un \textbf{master} e uno o più \textbf{slave}. Il master serve
        letture e scritture e replica ogni scrittura su ogni slave; gli slave sono \textbf{repliche in sola lettura}.
        \textit{Disponibilità}: se il master si guasta, il sistema può continuare in sola lettura, e uno slave può
        eventualmente essere promosso a master. Funziona bene quando ci sono \textbf{molte letture e poche scritture}.
  \item \textbf{Master-master.} Tutte le istanze sono master e servono sia letture sia scritture. A ogni scrittura i master
        si coordinano e si sincronizzano, mentre le letture sono gestite in modo indipendente. \textit{Disponibilità}: se un
        master si guasta, il sistema continua con l'altro. La sincronizzazione però introduce un \textbf{sovraccarico}.
\end{itemize}

\subsection{Federazione e sharding}

Le slide le elencano fra le tecniche di replica, ma a rigore sono tecniche di \textbf{partizionamento}: invece di copiare
gli stessi dati su più database, si \textbf{dividono} i dati fra più database.

\Needspace{12\baselineskip}
\noindent\begin{minipage}[c]{0.56\textwidth}
\textbf{Federazione.} Un unico grande database viene diviso in più database \textbf{per funzione}: invece di un database
monolitico, uno per i prodotti, uno per gli utenti, uno per gli ordini. Gli svantaggi:
\begin{itemize}
  \item non serve a nulla se lo schema ha un'unica tabella molto grande;
  \item è più difficile fare \textit{join} fra dati che stanno in database diversi;
  \item è più difficile far rispettare vincoli fra tabelle (ad esempio le chiavi esterne);
  \item può richiedere di denormalizzare.
\end{itemize}
\end{minipage}\hfill
\begin{minipage}[c]{0.4\textwidth}
\centering
\begin{tikzpicture}
  \node[servizio, fill=swlight, draw=swgray] (a) at (1.6,1.3) {applicazione};
  \node[db] (p) at (0,-0.4) {Prodotti};
  \node[db] (u) at (1.6,-0.4) {Utenti};
  \node[db] (o) at (3.2,-0.4) {Ordini};
  \foreach \d in {p,u,o} \draw[fr] (a) -- (\d);
\end{tikzpicture}
\end{minipage}

\Needspace{12\baselineskip}
\noindent\begin{minipage}[c]{0.56\textwidth}
\textbf{Sharding.} I dati vengono distribuiti su più database in modo che ciascuno gestisca solo un \textbf{sottoinsieme}
dei record della stessa tabella: ad esempio gli utenti con il nome dalla A alla M in un database, quelli dalla N alla Z in
un altro. Gli svantaggi:
\begin{itemize}
  \item la logica applicativa diventa più complessa, perché deve sapere dove sta ogni dato;
  \item gli \textit{shard} possono \textbf{sbilanciarsi} (in Italia i cognomi che iniziano per R sono molti di più di
        quelli che iniziano per Q);
  \item unire i dati è difficile, e far rispettare i vincoli lo è ancora di più.
\end{itemize}
\end{minipage}\hfill
\begin{minipage}[c]{0.4\textwidth}
\centering
\begin{tikzpicture}
  \node[servizio, fill=swlight, draw=swgray] (a) at (1.2,2.2) {applicazione};
  \node[rectangle, fill=swred!80!black, text=white, font=\scriptsize\bfseries, minimum height=0.6cm] (lb) at (1.2,1.0)
    {instradamento};
  \draw[fr] (a) -- node[font=\tiny, right]{utente «Luigi»} (lb);
  \node[db] (s1) at (0,-0.5) {Utenti\\ A--M};
  \node[db] (s2) at (2.4,-0.5) {Utenti\\ N--Z};
  \draw[fr] (lb) -- (s1); \draw[fr, line width=1.2pt, draw=swred] (lb) -- (s2);
\end{tikzpicture}
\end{minipage}

\subsection{Denormalizzazione}

La \term{denormalizzazione} migliora le prestazioni in lettura a scapito di quelle in scrittura (e, naturalmente, della
normalizzazione). Si scrivono \textbf{copie ridondanti} degli stessi dati in più tabelle, per evitare \textit{join} costosi.
Il prezzo è che mantenere le copie consistenti fra loro diventa un problema. Alcuni RDBMS, come PostgreSQL, offrono le
\textbf{viste materializzate}, che memorizzano il risultato di un'interrogazione e si occupano di gestire le copie
ridondanti (in PostgreSQL la vista va aggiornata esplicitamente con \code{REFRESH MATERIALIZED VIEW}; altri DBMS la
aggiornano automaticamente).

\subsection{I database NoSQL e la consistenza finale}

I database \term{NoSQL} sono più scalabili, a costo di rinunciare a vere transazioni ACID (e a parte delle capacità di
interrogazione). Le loro proprietà si riassumono spesso con l'acronimo \textbf{BASE}:
\begin{itemize}
  \item \textbf{Basically Available}: il sistema garantisce la disponibilità;
  \item \textbf{Soft state}: lo stato del sistema può cambiare nel tempo, anche senza input;
  \item \textbf{Eventual consistency}: il sistema diventa consistente \textbf{dopo un certo tempo}, a patto di non ricevere
        nuovi input in quel periodo.
\end{itemize}

\dfn{Consistenza finale (eventual consistency)}{
  Garanzia informale che, se un dato non riceve nuovi aggiornamenti, \textbf{prima o poi} tutti gli accessi a quel dato
  restituiranno l'ultimo valore scritto.
}

\begin{figure}[H]
\centering
\begin{tikzpicture}[
    nodo/.style={rectangle, rounded corners=2pt, draw=swblue!70!black, line width=0.7pt, minimum width=1.2cm,
                 minimum height=0.75cm, font=\scriptsize, align=center},
    vecchio/.style={nodo, fill=white},
    nuovo/.style={nodo, fill=swgreen!22},
    ev/.style={font=\scriptsize\bfseries, text=#1},
    passo/.style={chip=swgray, font=\scriptsize\bfseries}]
  \foreach \r/\stati/\tit in {
      0/{0,0,0,0,0}/{1. Tutti i nodi hanno X = 1},
      1/{0,1,0,0,0}/{2. Una scrittura X = 2 arriva al nodo 2},
      2/{1,1,1,0,1}/{3. La modifica si propaga, ma non ovunque},
      3/{1,1,1,1,1}/{4. Prima o poi tutti i nodi sono allineati}} {
    \node[passo, anchor=west] at (-0.6,-\r*1.55+0.55) {\tit};
    \foreach \v [count=\c from 1] in \stati {
      \ifnum\v=1
        \node[nuovo] (n\r\c) at (\c*1.55,-\r*1.55-0.15) {Nodo \c\\ \textbf{X = 2}};
      \else
        \node[vecchio] (n\r\c) at (\c*1.55,-\r*1.55-0.15) {Nodo \c\\ X = 1};
      \fi
    }
  }
  \draw[fr, draw=swred, line width=1pt] (2*1.55,-1*1.55+0.55) -- (n12.north);
  \node[ev=swred, anchor=west] at (9.0,-1*1.55-0.15) {\code{write X=2}};
  \node[ev=swgreen, anchor=west] at (9.0,-2*1.55+0.05) {\code{read X} sul nodo 1 $\rightarrow$ 2};
  \node[ev=swred, anchor=west] at (9.0,-2*1.55-0.35) {\code{read X} sul nodo 4 $\rightarrow$ \textbf{1}!};
  \node[ev=swgreen, anchor=west] at (9.0,-3*1.55-0.15) {ogni \code{read X} $\rightarrow$ 2};
\end{tikzpicture}
\caption{La consistenza finale in un cluster NoSQL di cinque nodi. Al passo 3 due letture contemporanee danno risultati
         diversi: è il prezzo della scalabilità.}
\end{figure}

Esistono diversi modelli di database NoSQL:
\begin{center}
\renewcommand{\arraystretch}{1.35}
\begin{tabular}{@{}>{\bfseries}l >{\raggedright\arraybackslash}p{8.6cm} l@{}}
  \toprule
  \normalfont\textbf{Modello} & \textbf{Come organizza i dati} & \textbf{Esempi} \\
  \midrule
  Chiave-valore & coppie chiave/valore, tipicamente in memoria & Redis, Memcached \\
  Documentali   & documenti (XML, JSON\ldots) che contengono tutte le informazioni di un oggetto & MongoDB \\
  A grafo       & ogni record è un nodo, le relazioni fra record sono archi & Neo4J \\
  Vettoriali    & memorizzazione e ricerca di grandi vettori, spesso prodotti da applicazioni di intelligenza
                  artificiale & Milvus \\
  \bottomrule
\end{tabular}
\end{center}

% ============================================================
\section{Elaborazione asincrona}

Man mano che un sistema cresce, può capitare di dover svolgere elaborazioni che non si possono fare \textbf{in linea}
mentre si serve una richiesta, perché introdurrebbero una latenza inaccettabile. Nell'esempio, a ogni richiesta che produce
un errore, l'applicazione chiede a un \code{ErrorHandler} di registrarlo, avvisare gli amministratori e analizzarlo, e solo
dopo risponde al client.

\begin{figure}[H]
\centering
\begin{tikzpicture}
  \foreach \x/\n in {0/Client, 2.0/Application, 4.2/ErrorHandler, 6.4/Logger, 8.4/Mailer, 10.6/SecurityAnalyzer}
    {\node[partec] at (\x,0) {\n}; \draw[vita] (\x,-0.3) -- (\x,-4.7);}
  \draw[sinc] (0,-0.8) -- node[msgl]{request(r)} (1.92,-0.8);
  \attiv{2}{-0.8}{-4.2}
  \draw[sinc] (2.08,-1.2) -- node[msgl]{handle(e)} (4.12,-1.2);
  \attiv{4.2}{-1.2}{-3.85}
  \draw[sinc] (4.28,-1.6) -- node[msgl]{logError(e)} (6.32,-1.6); \attiv{6.4}{-1.6}{-1.95}
  \draw[ritorno] (6.32,-1.95) -- (4.28,-1.95);
  \draw[sinc] (4.28,-2.3) -- node[msgl]{notifyAdmins(e)} (8.32,-2.3); \attiv{8.4}{-2.3}{-2.75}
  \draw[ritorno] (8.32,-2.75) -- (4.28,-2.75);
  \draw[sinc] (4.28,-3.1) -- node[msgl, pos=0.65]{analyze(e)} (10.52,-3.1); \attiv{10.6}{-3.1}{-3.65}
  \draw[ritorno] (10.52,-3.65) -- (4.28,-3.65);
  \draw[ritorno] (4.12,-3.85) -- (2.08,-3.85);
  \draw[ritorno] (1.92,-4.2) -- node[msgl]{response} (0,-4.2);
  \draw[swred, dashed, line width=0.9pt] (-0.6,-1.15) -- (11.5,-1.15);
  \draw[swred, dashed, line width=0.9pt] (-0.6,-3.9) -- (11.5,-3.9);
  \draw[{Stealth[length=2mm]}-{Stealth[length=2mm]}, swred, line width=0.9pt] (-0.5,-1.15) -- (-0.5,-3.9);
  \node[font=\scriptsize, text=swred, align=right, anchor=east] at (-0.6,-2.5) {latenza\\ inutile};
\end{tikzpicture}
\caption{Elaborazione sincrona: il client aspetta che l'errore sia registrato, notificato e analizzato, anche se a lui
         non serve.}
\end{figure}

\subsection{Code di messaggi e message broker}

L'asincronismo aiuta a ridurre i tempi di risposta. Un modo di realizzarlo sono le \term{code di messaggi}
(\textit{Message Queue}, MQ): componenti progettati per ricevere, conservare e consegnare messaggi. Quando un compito è
troppo lento per essere svolto in linea:
\begin{enumerate}
  \item l'applicazione \textbf{pubblica} un messaggio in una coda, con i dettagli del compito da svolgere;
  \item un \textbf{worker}, in modo asincrono e magari più tardi, fa da \textbf{consumatore}: preleva il messaggio dalla coda e
        svolge davvero il compito;
  \item l'utente \textbf{non resta bloccato} ad aspettare la fine del compito lungo, che viene completato in background
        appena uno dei worker (che possono essere molti) è disponibile.
\end{enumerate}

\begin{figure}[H]
\centering
\begin{tikzpicture}[
    prod/.style={rectangle, rounded corners=3pt, fill=swteal!80!black, text=white, font=\scriptsize\bfseries,
                 minimum width=1.9cm, minimum height=0.6cm},
    cons/.style={rectangle, rounded corners=3pt, fill=swgreen!75!black, text=white, font=\scriptsize\bfseries,
                 minimum width=1.9cm, minimum height=0.55cm},
    coda/.style={cylinder, shape border rotate=0, aspect=0.3, draw=#1!70!black, fill=#1!18, line width=0.8pt,
                 minimum width=1.0cm, minimum height=2.7cm, font=\scriptsize\bfseries}]
  \draw[rounded corners=4pt, fill=yellow!35, draw=q5!60!black, line width=0.8pt] (2.5,1.5) rectangle (6.7,-2.6);
  \node[font=\small\bfseries] at (4.6,1.15) {Message broker};
  \node[coda=q4] (q1) at (4.6,0.2) {MQ 1 \quad\imail\,\imail};
  \node[coda=q7] (q2) at (4.6,-1.6) {MQ 2 \quad\imail\,\imail};
  \node[prod] (p1) at (0,0.6) {Producer 1};
  \node[prod] (p2) at (0,-0.4) {Producer 2};
  \node[prod] (p3) at (0,-1.4) {Producer 3};
  \node[cons] (c1) at (9.2,1.0) {Consumer 1};
  \node[cons] (c2) at (9.2,0.2) {Consumer 2};
  \node[cons] (c3) at (9.2,-0.9) {Consumer 3};
  \node[cons] (c4) at (9.2,-2.0) {Consumer $n$};
  \draw[fr] (p1.east) -- (q1.west); \draw[fr] (p2.east) -- (q1.west); \draw[fr] (p2.east) -- (q2.west);
  \draw[fr] (p3.east) -- (q2.west);
  \draw[fr] (q1.east) -- (c1.west); \draw[fr] (q1.east) -- (c2.west);
  \draw[fr] (q2.east) -- (c3.west); \draw[fr] (q2.east) -- (c4.west);
\end{tikzpicture}
\caption{Produttori e consumatori non si conoscono: comunicano solo attraverso le code del broker.}
\end{figure}

\begin{figure}[H]
\centering
\begin{tikzpicture}
  \foreach \x/\n in {0/Client, 1.9/Application, 3.9/MessageQueue, 5.9/ErrorHandler, 7.7/Logger, 9.2/Mailer,
                     11.0/SecurityAnalyzer}
    {\node[partec] at (\x,0) {\n}; \draw[vita] (\x,-0.3) -- (\x,-5.0);}
  \draw[sinc] (0,-0.8) -- node[msgl]{request(r)} (1.82,-0.8);
  \attiv{1.9}{-0.8}{-1.7}
  \draw[asinc] (1.98,-1.15) -- node[msgl]{store(e)} (3.82,-1.15);
  \attiv{3.9}{-1.15}{-1.5}
  \draw[ritorno] (1.82,-1.7) -- node[msgl]{response} (0,-1.7);
  \attiv{5.9}{-0.6}{-4.6}
  \draw[sinc] (5.82,-2.2) -- node[msgl]{dequeue()} (3.98,-2.2);
  \attiv{3.9}{-2.2}{-2.55}
  \draw[ritorno] (3.98,-2.55) -- node[msgl]{e} (5.82,-2.55);
  \draw[sinc] (5.98,-3.0) -- node[msgl]{logError(e)} (7.62,-3.0); \attiv{7.7}{-3.0}{-3.3}
  \draw[sinc] (5.98,-3.6) -- node[msgl]{notifyAdmins(e)} (9.12,-3.6); \attiv{9.2}{-3.6}{-3.9}
  \draw[sinc] (5.98,-4.2) -- node[msgl, pos=0.7]{analyze(e)} (10.92,-4.2); \attiv{11.0}{-4.2}{-4.5}
  \draw[{Stealth[length=2mm]}-{Stealth[length=2mm]}, swgreen, line width=0.9pt] (-0.5,-0.8) -- (-0.5,-1.7);
  \node[font=\scriptsize, text=swgreen!70!black, align=right, anchor=east] at (-0.6,-1.25) {latenza\\ ridotta};
  \draw[swred, dashed, line width=0.8pt, rounded corners=3pt] (5.4,-1.95) rectangle (11.9,-4.75);
  \node[font=\scriptsize, text=swred, align=left, anchor=west] at (12.0,-3.35) {gestione\\ dell'errore\\ asincrona};
\end{tikzpicture}
\caption{Con una coda di messaggi l'applicazione deposita l'errore e risponde subito; l'\code{ErrorHandler} lo preleva e lo
         gestisce per conto suo.}
\end{figure}

Fra i \term{message broker} open source più diffusi ci sono RabbitMQ, Apache ActiveMQ e BlazingMQ.

\warning{{I message broker non sono oggetti banali}}{
  Internamente sono organizzati in architetture distribuite complesse; l'instradamento dei messaggi può essere più
  sofisticato di una semplice interrogazione da parte dei consumatori; e broker e configurazioni diversi garantiscono
  \textbf{proprietà diverse}: la consegna di un messaggio può essere garantita o no, e lo stesso messaggio può essere
  consegnato a \textbf{più di un} consumatore. Bisogna valutare con attenzione questi aspetti rispetto al problema da
  risolvere: un messaggio consegnato due volte, ad esempio, non deve far addebitare due volte un pagamento.
}

\subsection{Un altro scenario: lo stesso utente su più dispositivi}

Le code non servono solo a spostare in background i compiti lenti. Un altro scenario tipico è quello delle applicazioni a
cui lo \textbf{stesso utente} accede da \textbf{più dispositivi}, anche contemporaneamente: Microsoft Teams, per esempio,
può essere aperto nello stesso momento sul telefono, sul portatile e nel browser. Quando arriva un nuovo messaggio deve
comparire su tutti i dispositivi; quando l'utente lo legge su uno, gli altri devono smettere di segnalarlo come non letto.

Se fosse il servizio di chat a chiamare i dispositivi uno per uno, dovrebbe sapere in ogni momento quali sono connessi,
aspettare la risposta di ciascuno e gestire quelli momentaneamente irraggiungibili (un telefono in galleria). Con un
message broker, invece, il servizio \textbf{pubblica} l'evento \textbf{una sola volta}; ogni dispositivo connesso è un
\textbf{consumatore} iscritto agli eventi del proprio utente, e il broker consegna a ciascuno la sua copia. È proprio il
caso, citato nel riquadro precedente, in cui lo stesso messaggio arriva a \textbf{più} consumatori: qui è esattamente
ciò che si vuole. E se il broker conserva i messaggi per i consumatori disconnessi, un dispositivo che si riconnette trova
in coda gli eventi che si era perso.

\Needspace{12\baselineskip}
\begin{figure}[H]
\centering
\begin{tikzpicture}[
    att/.style={rectangle, rounded corners=3pt, fill=swteal!80!black, text=white, font=\scriptsize\bfseries,
                minimum width=2.2cm, minimum height=0.6cm, align=center},
    disp/.style={rectangle, rounded corners=3pt, fill=swgreen!75!black, text=white, font=\scriptsize\bfseries,
                 minimum width=2.6cm, minimum height=0.55cm},
    lab/.style={font=\tiny, above, inner sep=1.5pt, align=center}]
  \node[att] (b) at (-0.6,0) {app di Bruno};
  \node[servizio] (s) at (3.4,0) {servizio\\ di chat};
  \draw[rounded corners=4pt, fill=yellow!35, draw=q5!60!black, line width=0.8pt] (5.4,1.25) rectangle (8.6,-1.35);
  \node[font=\small\bfseries] at (7.0,0.95) {Message broker};
  \node[cylinder, shape border rotate=0, aspect=0.3, draw=q4!70!black, fill=q4!18, line width=0.8pt,
        minimum width=0.9cm, minimum height=2.1cm, font=\scriptsize\bfseries] (q) at (7.0,-0.2)
        {eventi di Anna \quad\imail};
  \node[disp] (d1) at (11.4,0.9)  {telefono di Anna};
  \node[disp] (d2) at (11.4,0.0)  {portatile di Anna};
  \node[disp] (d3) at (11.4,-0.9) {Teams web di Anna};
  \draw[fr] (b) -- node[lab]{nuovo messaggio\\ per Anna} (s);
  \draw[fr] (s.east) -- node[lab]{pubblica\\ (una volta)} (q.west |- s.east);
  \draw[fr] (q.east) -- (d1.west);
  \draw[fr] (q.east) -- (d2.west);
  \draw[fr] (q.east) -- (d3.west);
  \node[font=\tiny, text=swgreen!60!black, align=center] at (11.4,-1.55) {una copia a ogni dispositivo connesso};
\end{tikzpicture}
\caption{Lo stesso utente collegato da tre dispositivi: il servizio pubblica l'evento una volta sola e il broker lo
         consegna a ciascun dispositivo.}
\end{figure}

% ============================================================
\section{Microservizi}

\Needspace{16\baselineskip}
\noindent\begin{minipage}[t]{0.49\textwidth}
\gbox{Architettura monolitica}{q6}{\raggedright
\begin{itemize}[leftmargin=1.1em]
  \item componenti e moduli \textbf{strettamente impacchettati} in un'unica unità;
  \item un'unica \textbf{codebase};
  \item un'unica unità di \textbf{deployment}, autosufficiente;
  \item le responsabilità di business sono \textbf{accoppiate}.
\end{itemize}
\centering
\begin{tikzpicture}
  \draw[rounded corners=3pt, fill=q6!10, draw=q6!70!black, line width=0.9pt] (0,0) rectangle (4.2,1.4);
  \foreach \x/\n in {0.55/Utenti, 1.6/Ordini, 2.65/Catalogo, 3.65/Pagam.}
    \node[font=\tiny\bfseries, rounded corners=1pt, draw=q6!50, fill=white, minimum width=0.9cm, minimum height=0.45cm]
      at (\x,0.7) {\n};
  \node[font=\tiny, text=q6!80!black] at (2.1,0.18) {un solo processo, un solo deployment};
\end{tikzpicture}}
\end{minipage}\hfill
\begin{minipage}[t]{0.49\textwidth}
\gbox{Architettura a microservizi}{swgreen}{\raggedright
\begin{itemize}[leftmargin=1.1em]
  \item il software è strutturato come un insieme di \textbf{servizi};
  \item ogni servizio ha la \textbf{propria codebase}, anche in linguaggi diversi (\textit{poliglotta});
  \item ogni servizio si può \textbf{rilasciare in modo indipendente};
  \item i servizi sono organizzati attorno alle \textbf{capacità di business};
  \item i servizi sono \textbf{debolmente accoppiati}.
\end{itemize}
\centering
\begin{tikzpicture}
  \foreach \x/\y/\n/\k in {0.5/1.0/Utenti/a, 2.1/1.0/Ordini/b, 3.7/1.0/Catalogo/c, 2.1/0/Pagam./d}
    \node[servizio, minimum width=1.2cm, minimum height=0.5cm, font=\tiny\bfseries] (\k) at (\x,\y) {\n};
  \draw[frd] (a) -- (b); \draw[frd] (b) -- (c); \draw[frd] (b) -- (d);
\end{tikzpicture}}
\end{minipage}

\subsection{Perché mai dovrei farlo?}

All'inizio, o per software poco complesso, il monolite ha vantaggi concreti:
\begin{itemize}
  \item \textbf{deployment facile}: c'è un solo artefatto da costruire (ad esempio un unico file WAR da caricare su un
        application server come Apache Tomcat, che parla con un unico database);
  \item \textbf{sviluppo facile}: un'unica codebase, e una gestione dei dati semplice.
\end{itemize}

I problemi arrivano quando il monolite \textbf{cresce in complessità}. Le slide lo paragonano a Sisifo, condannato a
spingere in eterno un masso su per un monte:
\begin{itemize}
  \item l'architettura interna diventa sempre più difficile da mantenere e da far evolvere;
  \item i nuovi rilasci richiedono sempre più tempo (mesi), e aggiungere funzionalità diventa lento;
  \item il ciclo di \textit{build}, test e rilascio diventa lungo e complicato: \textit{chi ha rotto la build?};
  \item il deployment diventa sempre più difficile: \textit{di chi è il modulo che non funziona?};
  \item l'innovazione si ferma.
\end{itemize}

\subsection{Le motivazioni: consegnare spesso e bene}

Il mondo di oggi cambia rapidamente, le aziende devono essere agili e innovare in fretta, e il software deve essere
consegnato \textbf{rapidamente, frequentemente e in modo affidabile}. Le quattro metriche \textbf{DORA}
(\textit{DevOps Research and Assessment}) misurano proprio questo:

\begin{center}
\renewcommand{\arraystretch}{1.35}
\begin{tabular}{@{}>{\bfseries}l >{\raggedright\arraybackslash}p{9.6cm}@{}}
  \toprule
  \normalfont\textbf{Metrica} & \textbf{Che cosa misura} \\
  \midrule
  Deployment Frequency     & quanto spesso si rilascia in produzione \\
  Lead Time for Changes    & quanto tempo impiega un \textit{commit} ad arrivare in produzione \\
  Change Failure Rate      & la percentuale di rilasci che causano un guasto in produzione \\
  Time to Restore Service  & quanto tempo serve per riprendersi da un guasto in produzione \\
  \bottomrule
\end{tabular}
\end{center}

Dividere il monolite in microservizi aiuta a rendere gestibili i sistemi grandi:
\begin{itemize}
  \item ogni microservizio è molto più piccolo del sistema intero, quindi più facile da capire, mantenere, far evolvere e
        testare;
  \item ogni microservizio si può costruire e rilasciare in modo indipendente;
  \item una nuova funzionalità probabilmente tocca un solo microservizio;
  \item i rilasci sono più frequenti: per aggiungere una funzionalità basta rilasciare di nuovo un solo microservizio.
\end{itemize}

\subsection{Scalare un monolite e scalare dei microservizi}

\begin{figure}[H]
\centering
\begin{tikzpicture}[
    srv/.style={rectangle, rounded corners=3pt, draw=swteal!80!black, fill=swteal!8, line width=0.8pt,
                minimum width=2.0cm, minimum height=1.7cm},
    f/.style={rectangle, rounded corners=1pt, draw=#1!70!black, fill=#1!25, line width=0.5pt,
              minimum width=0.42cm, minimum height=0.32cm, inner sep=0pt}]
  % monolite replicato
  \node[font=\small\bfseries] at (3.3,1.4) {Monolite};
  \foreach \s in {0,1,2} {
    \node[srv] at (\s*2.2+1.1,0) {};
    \draw[rounded corners=2pt, fill=q6!10, draw=q6!60] (\s*2.2+0.4,-0.65) rectangle (\s*2.2+1.8,0.65);
    \foreach \c/\x/\y in {q4/0.65/0.35, q5/1.1/0.35, q7/1.55/0.35, q8/0.65/-0.05, swblue/1.1/-0.05, q9/1.55/-0.05}
      \node[f=\c] at (\s*2.2+\x,\y) {};
  }
  \node[note, align=center] at (3.3,-1.2) {si replica \textbf{tutto}, anche le funzioni\\ che non ne hanno bisogno};
  % microservizi distribuiti
  \begin{scope}[shift={(8.0,0)}]
    \node[font=\small\bfseries] at (3.3,1.4) {Microservizi};
    \foreach \s in {0,1,2} \node[srv] at (\s*2.2+1.1,0) {};
    % server 1: q4, q5, q7 ; server 2: q4, q4, q8 ; server 3: q4, swblue, q9
    \foreach \c/\x/\y in {q4/0.7/0.35, q5/1.5/0.35, q7/0.7/-0.3, q4/1.5/-0.3} \node[f=\c, minimum width=0.6cm, minimum height=0.45cm] at (\x,\y) {};
    \foreach \c/\x/\y in {q4/0.7/0.35, q8/1.5/0.35, q4/0.7/-0.3, swblue/1.5/-0.3} \node[f=\c, minimum width=0.6cm, minimum height=0.45cm] at (2.2+\x,\y) {};
    \foreach \c/\x/\y in {q4/0.7/0.35, q9/1.5/0.35, q4/0.7/-0.3, q8/1.5/-0.3} \node[f=\c, minimum width=0.6cm, minimum height=0.45cm] at (4.4+\x,\y) {};
    \node[note, align=center] at (3.3,-1.2) {si replica \textbf{solo} ciò che serve\\ (qui il servizio rosso, il più richiesto)};
  \end{scope}
\end{tikzpicture}
\caption{Il monolite mette tutte le funzionalità in un unico processo e scala replicandolo per intero; i microservizi
         mettono ogni funzionalità in un processo distinto e scalano distribuendo i servizi sui server, replicando quelli che
         servono.}
\end{figure}

\subsection{Conviene davvero?}

\Needspace{18\baselineskip}
\noindent\begin{minipage}[c]{0.5\textwidth}
\textbf{Non prendete neanche in considerazione i microservizi, a meno che il sistema non sia complesso.} La complessità che
spinge verso i microservizi può venire da:
\begin{itemize}
  \item le dimensioni pure, con team di sviluppo molto grandi;
  \item la \textbf{multi-tenancy} (lo stesso sistema serve molti clienti separati);
  \item il dover supportare modelli di interazione diversi (web, app, API per terzi\ldots);
  \item la necessità di scalare.
\end{itemize}
Il grafico (Fowler, \textit{Microservice Premium}) lo mostra: per sistemi poco complessi, il sovraccarico dei microservizi
\textbf{riduce} la produttività. Quando la complessità cresce, la produttività del monolite crolla; quella dei microservizi
cala molto meno, grazie al minore accoppiamento.
\end{minipage}\hfill
\begin{minipage}[c]{0.46\textwidth}
\centering
\begin{tikzpicture}[x=0.75cm, y=0.75cm]
  \draw[-{Stealth[length=2mm]}, line width=0.8pt] (0,0) -- (8.4,0) node[below left, note]{complessità};
  \draw[-{Stealth[length=2mm]}, line width=0.8pt] (0,0) -- (0,5.6) node[above right, note]{produttività};
  \draw[line width=1.3pt, draw=swblue, domain=0:8.1, samples=60] plot (\x, {1.0+3.4/(1+exp(1.6*(\x-3.2)))});
  \draw[line width=1.3pt, draw=sworange, domain=0:8.1, samples=60] plot (\x, {2.3+0.7/(1+exp(1.4*(\x-5.4)))});
  \node[note, text=swblue, font=\scriptsize\bfseries, anchor=west] at (0.2,4.85) {monolite};
  \node[note, text=sworange!90!black, font=\scriptsize\bfseries, anchor=west] at (5.9,2.75) {microservizi};
  \draw[{Stealth[length=1.8mm]}-{Stealth[length=1.8mm]}, swred, line width=0.8pt] (0.9,3.05) -- (0.9,4.3);
  \node[note, font=\tiny, text=swred, anchor=west, align=left] at (1.0,3.6) {costo fisso dei\\ microservizi};
  \draw[{Stealth[length=1.8mm]}-{Stealth[length=1.8mm]}, swgreen, line width=0.8pt] (6.4,1.05) -- (6.4,2.35);
  \node[note, font=\tiny, text=swgreen!70!black, anchor=west, align=left] at (6.5,1.6) {guadagno a\\ complessità alta};
\end{tikzpicture}
\end{minipage}

Adottare i microservizi significa rispondere a quattro domande: come progettare i microservizi, come gestire la
comunicazione fra servizi, come gestire i dati, come organizzare i team di sviluppo. Le vediamo una alla volta.

\subsection{Progettare i microservizi}

Definire i \textbf{confini} di ogni servizio è una delle sfide più grandi. La regola generale è che un servizio dovrebbe
fare «\textbf{una cosa sola}» (è il Single Responsibility Principle del Capitolo 10, applicato ai servizi), ma nessun
procedimento meccanico garantisce il risultato giusto. Viene in aiuto il \term{Domain-Driven Design} (DDD), un approccio alla
progettazione centrato sul modello del dominio. Le slide lo chiamano \textit{Domain-Driven Development}, ma il nome
corretto, introdotto da Eric Evans nel 2003, è \textit{Domain-Driven \textbf{Design}}. Il procedimento ha quattro passi.

\begin{center}
\begin{tikzpicture}[
    fase/.style={phase=#1, minimum width=3.0cm, minimum height=0.95cm, font=\scriptsize\bfseries}]
  \node[fase=ph1] (a) at (0,0)   {Analizzare\\ il dominio};
  \node[fase=ph2] (b) at (3.5,0) {Definire i\\ bounded context};
  \node[fase=ph3] (c) at (7.0,0) {Definire entità,\\ aggregati, servizi};
  \node[fase=ph4] (d) at (10.5,0) {Identificare i\\ microservizi};
  \draw[flow] (a) -- (b); \draw[flow] (b) -- (c); \draw[flow] (c) -- (d);
\end{tikzpicture}
\end{center}

\begin{enumerate}
  \item \textbf{Analizzare il dominio.} Si comprendono i requisiti funzionali; il risultato è una descrizione informale del
        dominio del sistema, che si può raffinare in un insieme di modelli di dominio (Capitolo 7).
  \item \textbf{Definire i bounded context.} Ogni \term{bounded context} (contesto delimitato) contiene un modello di dominio
        che rappresenta un particolare \textbf{sottodominio}. Unificare tutto il dominio in un solo modello potrebbe non essere
        né possibile né conveniente.
  \item \textbf{Definire entità, aggregati e servizi di dominio.} Per ogni bounded context si applicano i pattern del DDD:
        \begin{itemize}
          \item le \textbf{entità}: oggetti con un'\textbf{identità} che persiste nel tempo (un cliente resta lo stesso
                cliente anche se cambia indirizzo);
          \item gli \textbf{aggregati}: confini di \textbf{consistenza} attorno a una o più entità, che modellano gli
                invarianti da rispettare in una transazione (un ordine con le sue righe: il totale deve sempre essere la
                somma delle righe);
          \item i \textbf{servizi di dominio}: oggetti che implementano logica senza conservare uno stato.
        \end{itemize}
  \item \textbf{Identificare i microservizi.} Si parte da un bounded context: le funzionalità di un microservizio non
        dovrebbero estendersi su più bounded context. Aggregati e servizi di dominio sono buoni candidati a diventare
        microservizi.
\end{enumerate}

\begin{esempio}{La stessa parola, due contesti}
In un e-commerce, «prodotto» significa cose diverse a seconda di chi ne parla.

\begin{center}
\begin{tikzpicture}
  \node[draw=ph2!70!black, fill=ph2!6, rounded corners=6pt, line width=0.9pt, minimum width=5.4cm, minimum height=2.6cm]
    (cat) at (0,0) {};
  \node[font=\scriptsize\bfseries, text=ph2!80!black, anchor=north] at (cat.north) {Bounded context: Catalogo};
  \umlclass[3.4cm]{p1}{(0,-0.25)}{Prodotto}{nome, descrizione\\ foto, prezzo di listino}{}
  \node[draw=ph4!70!black, fill=ph4!6, rounded corners=6pt, line width=0.9pt, minimum width=5.4cm, minimum height=2.6cm]
    (mag) at (7.0,0) {};
  \node[font=\scriptsize\bfseries, text=ph4!80!black, anchor=north] at (mag.north) {Bounded context: Magazzino};
  \umlclass[3.4cm]{p2}{(7.0,-0.25)}{Prodotto}{codice, quantità disponibile\\ scaffale, peso}{}
  \draw[dashed, line width=0.7pt, draw=swgray] (p1.east) -- node[note, above]{stesso codice} (p2.west);
\end{tikzpicture}
\end{center}
Forzare un unico modello «Prodotto» per tutto il sistema produrrebbe una classe enorme, che cambia ogni volta che cambia
il catalogo \textit{o} il magazzino: è la God Class del Capitolo 10, a livello di sistema. Con due bounded context, ognuno
ha il proprio modello, e i due condividono solo l'identificativo. Catalogo e Magazzino sono due buoni candidati a diventare
due microservizi.
\end{esempio}

\subsection{La comunicazione fra microservizi}

In un monolite i componenti comunicano con \textbf{chiamate di metodo}. In un'architettura a microservizi la comunicazione
dovrebbe essere \textbf{indipendente dalla tecnologia}: spesso i team usano i protocolli e le tecnologie su cui è costruito
il Web. La comunicazione può essere:
\begin{itemize}
  \item \textbf{sincrona}: REST (API HTTP orientate alle risorse), Apache Thrift (che supporta più protocolli: binari,
        basati su JSON\ldots);
  \item \textbf{asincrona}: code di messaggi.
\end{itemize}

\begin{figure}[H]
\centering
\begin{tikzpicture}
  % sincrona
  \node[font=\small\bfseries] at (2.3,0.7) {Comunicazione sincrona};
  \foreach \x/\n in {0/Frontend, 2.3/DataService, 4.6/LogService} {\node[partec] at (\x,0) {\n}; \draw[vita] (\x,-0.3) -- (\x,-3.2);}
  \draw[sinc] (-0.9,-0.7) -- node[msgl]{click()} (-0.08,-0.7);
  \attiv{0}{-0.7}{-2.9}
  \draw[sinc] (0.08,-1.1) -- node[msgl]{update(data)} (2.22,-1.1);
  \attiv{2.3}{-1.1}{-2.6}
  \draw[sinc] (2.38,-1.5) -- node[msgl]{log("update")} (4.52,-1.5);
  \attiv{4.6}{-1.5}{-2.2}
  \draw[ritorno] (4.52,-2.2) -- (2.38,-2.2);
  \draw[ritorno] (2.22,-2.6) -- (0.08,-2.6);
  % asincrona
  \begin{scope}[shift={(7.6,0)}]
    \node[font=\small\bfseries] at (3.2,0.7) {Comunicazione asincrona};
    \foreach \x/\n in {0/Frontend, 2.1/DataService, 4.2/Queue, 6.3/LogService} {\node[partec] at (\x,0) {\n}; \draw[vita] (\x,-0.3) -- (\x,-3.2);}
    \draw[sinc] (-0.9,-0.7) -- node[msgl]{click()} (-0.08,-0.7);
    \attiv{0}{-0.7}{-1.9}
    \draw[sinc] (0.08,-1.1) -- node[msgl]{update(data)} (2.02,-1.1);
    \attiv{2.1}{-1.1}{-1.75}
    \draw[asinc] (2.18,-1.45) -- node[msgl]{add("update")} (4.12,-1.45);
    \draw[ritorno] (2.02,-1.75) -- (0.08,-1.75);
    \attiv{6.3}{-2.2}{-3.0}
    \draw[sinc] (6.22,-2.4) -- node[msgl]{get()} (4.28,-2.4);
    \draw[ritorno] (4.28,-2.8) -- node[msgl]{"update"} (6.22,-2.8);
  \end{scope}
\end{tikzpicture}
\caption{A sinistra il \code{DataService} aspetta il \code{LogService} prima di rispondere; a destra deposita il messaggio
         nella coda e risponde subito, e il \code{LogService} lo preleva quando vuole.}
\end{figure}

Le chiamate sincrone possono penalizzare le prestazioni, e soprattutto, quando sono molte, si manifesta un \textbf{effetto
moltiplicativo} dei fermi: se per servire una richiesta devono funzionare tutti i servizi di una catena, la disponibilità
complessiva è il \textbf{prodotto} delle disponibilità dei singoli servizi. O si rende asincrona la comunicazione, o si
gestiscono i fermi.

\warning{{Si moltiplicano le disponibilità, non i fermi}}{
  Le slide dicono che il fermo del sistema diventa «il prodotto dei fermi dei singoli servizi». A moltiplicarsi, in realtà,
  sono le \textbf{disponibilità}: se una richiesta attraversa in modo sincrono $n$ servizi indipendenti, con disponibilità
  $A_1, \ldots, A_n$, la disponibilità complessiva è $A = A_1 \cdot A_2 \cdots A_n$. Con dieci servizi al 99{,}9\% ciascuno,
  $A = 0{,}999^{10} \approx 99{,}0\%$: si passa da circa 9 ore a circa 3 giorni e mezzo di fermo all'anno. Il fermo
  complessivo è quindi circa la \textbf{somma} dei fermi, non il prodotto.
}

\subsection{La gestione dei dati}

I monoliti sono tipicamente associati a un unico database \textbf{centralizzato}, spesso anche per via dei costi di licenza
dei DBMS. Anche i microservizi possono lavorare con un database centralizzato, ma allora le modifiche allo schema vanno
\textbf{coordinate} fra i team, e i servizi tornano ad essere accoppiati attraverso i dati.

L'alternativa è un \textbf{database per ogni servizio}: ogni microservizio (e ogni team) è davvero indipendente, e si può
scegliere per ognuno il tipo di database più adatto (\textit{polyglot persistence}: un database relazionale per gli ordini,
uno a documenti per il catalogo, uno chiave-valore per le sessioni\ldots).

\begin{figure}[H]
\centering
\begin{tikzpicture}
  % monolite
  \node[servizio, fill=q6!10, draw=q6!70!black, minimum height=1.3cm, minimum width=1.6cm] (m) at (0,0.4) {Monolite};
  \node[db] (md) at (0,-1.4) {DB};
  \draw[frd] (m) -- (md);
  \node[note, font=\scriptsize\bfseries] at (0,-2.4) {Monolite};
  % microservizi, database condiviso
  \begin{scope}[shift={(4.6,0)}]
    \node[servizio] (a) at (-1.2,0.9) {A}; \node[servizio] (b) at (0,0.9) {B}; \node[servizio] (c) at (1.2,0.9) {C};
    \node[db, minimum width=2.6cm] (sd) at (0,-1.4) {DB condiviso};
    \foreach \s in {a,b,c} \draw[frd] (\s) -- (sd);
    \node[note, font=\scriptsize\bfseries, align=center] at (0,-2.4) {Microservizi,\\ database condiviso};
  \end{scope}
  % un database per servizio
  \begin{scope}[shift={(10.0,0)}]
    \foreach \s/\x/\t in {a/-1.6/relazionale, b/0/documenti, c/1.6/chiave-valore} {
      \node[servizio] (\s) at (\x,0.9) {\MakeUppercase{\s}};
      \node[db, font=\tiny\bfseries] (\s-d) at (\x,-1.1) {DB\\ \t};
      \draw[frd] (\s) -- (\s-d);
    }
    \draw[frd, dashed] (a) -- (b); \draw[frd, dashed] (b) -- (c);
    \node[note, font=\scriptsize\bfseries, align=center] at (0,-2.4) {Microservizi,\\ un database per servizio};
  \end{scope}
\end{tikzpicture}
\caption{Tre modi di gestire i dati. Nel terzo i servizi si scambiano dati solo attraverso le loro interfacce (tratteggio).}
\end{figure}

Con i dati distribuiti su molti database, garantire una \textbf{consistenza forte} diventa sempre più difficile: spesso ci si
accontenta della \textbf{consistenza finale} vista per i database NoSQL. A un certo punto nel futuro, tutte le letture di
un'entità restituiranno il valore aggiornato.

\subsection{È una questione di compromessi}

\Needspace{10\baselineskip}
\begin{center}
\renewcommand{\arraystretch}{1.4}
\begin{tabular}{@{}>{\raggedright\arraybackslash}p{6.6cm} >{\raggedright\arraybackslash}p{6.6cm}@{}}
  \toprule
  \textcolor{swgreen}{\textbf{I microservizi danno\ldots}} & \textcolor{swred}{\textbf{\ldots ma non gratis}} \\
  \midrule
  confini forti fra i moduli & più complessità (un sistema distribuito) \\
  rilasci indipendenti & consistenza solo finale \\
  libertà di scelta delle tecnologie & più complessità operativa (molti servizi da rilasciare e monitorare) \\
  \bottomrule
\end{tabular}
\end{center}

Bisogna quindi avere una \textbf{buona ragione} per adottarli. Un caso spesso citato è quello di Istio, un noto progetto
open source che nel 2020 è \textbf{tornato} da un'architettura a microservizi a una monolitica, perché la complessità
operativa superava i benefici (Mendonça et al., \textit{The Monolith Strikes Back}, IEEE Software, 2021).

% ============================================================
\section{Un caso reale: un'applicazione a microservizi}

Per vedere che cosa succede quando questi concetti incontrano la pratica, le slide presentano un'esperienza recente: lo
sviluppo di un'applicazione complessa a microservizi nell'ambito del corso di \textit{Software Architecture Design}.

\subsection{Il progetto}

\begin{itemize}
  \item L'applicazione è stata sviluppata nell'arco di \textbf{tre edizioni successive} del corso. Nella prima, i team
        hanno implementato i singoli microservizi, \textbf{separatamente}; nella seconda e nella terza i servizi sono stati
        integrati e arricchiti di nuove funzionalità. Infine alcuni ricercatori a contratto l'hanno messa in produzione e
        hanno risolto diversi problemi di prestazioni e di sicurezza.
  \item Il prodotto finale è un insieme di \textbf{immagini di container} da eseguire su \textbf{Docker}.
  \item L'applicazione è un \textbf{gioco} per imparare la \textbf{copertura del codice} (\textit{code coverage}). Dai nomi
        dei servizi e delle API se ne intuisce la meccanica: lo studente scrive casi di test per una classe sotto test
        (\textit{class under test}), il sistema li compila e ne misura la copertura, e il risultato è confrontato con
        quello di avversari automatici, i \textit{robot}, i cui test sono generati con strumenti come EvoSuite.
\end{itemize}

\subsection{L'architettura}

Il diagramma di deployment mostra la struttura concettuale. Il client è un browser, che comunica con il server in HTTP
sulla porta 80. Sul server, sopra il sistema operativo, gira Docker, che fa da ambiente di esecuzione: ogni microservizio
è un \textbf{container}. Le sigle T1, \ldots, T8 indicano il \textbf{team} che ha sviluppato ciascun servizio (T23 è il
servizio nato dall'unione del lavoro dei team 2 e 3).

\Needspace{24\baselineskip}
\begin{figure}[H]
\centering
\begin{tikzpicture}
  % client
  \nodo{cli}{(-0.6,-0.6)}{2.4}{1.8}{«device»\\ \textbf{Client}}
  \node[compo, minimum width=1.7cm, minimum height=0.6cm, font=\tiny, inner sep=2pt] at (0.6,-0.15)
    {«component»\\ \textbf{Browser}};
  % server, sistema operativo, Docker
  \nodo{srv}{(4.0,-4.95)}{8.6}{9.1}{«device»\\ \textbf{Server}}
  \draw[line width=0.7pt, draw=swred!60!black, fill=swred!10] (4.3,-4.7) rectangle (12.3,3.3);
  \node[font=\tiny\bfseries] at (8.3,3.08) {«OS»};
  \draw[line width=0.7pt, draw=swblue!60!black, fill=swblue!8] (4.6,-4.45) rectangle (12.0,2.85);
  \node[font=\tiny, align=center] at (8.3,2.5) {«execution environment»\\ \textbf{Docker}};
  % i container
  \foreach \x/\y/\n/\c in {4.8/0.95/API Gateway/q7, 8.7/0.95/UI Gateway/q7,
                           4.8/-0.25/T1 -- Admin Service/q5, 8.7/-0.25/T23 -- Student Service/q5,
                           4.8/-1.45/T4 -- Game Repository/q5, 8.7/-1.45/T5 -- Game Front End/q5,
                           4.8/-2.65/T6 -- Game Engine/q5, 8.7/-2.65/T7 -- Student Test Runner/q5,
                           6.75/-3.85/T8 -- EvoSuite Runner/q5}
    \nodocolore{c}{(\x,\y)}{3.1}{0.85}{«container»\\ \textbf{\n}}{\c};
  % comunicazione
  \draw[line width=0.8pt] (1.8,0.35) -- node[font=\tiny, above, align=center]{«protocol»\\ HTTP : 80} (4.0,0.35);
\end{tikzpicture}
\caption{Il diagramma di deployment dell'applicazione (ridisegnato dalle slide): un container Docker per ogni
         microservizio, più due gateway.}
\end{figure}

Il diagramma dei componenti mostra chi chiama chi. Due componenti fanno da \textbf{gateway}: l'\textit{UI Gateway} porta le
richieste delle pagine ai tre servizi che producono l'interfaccia (il front end del gioco, il servizio degli studenti e
quello di amministrazione), mentre l'\textit{API Gateway} instrada le richieste di compilazione e di calcolo della copertura.
Molte chiamate, però, vanno \textbf{direttamente} da un servizio all'altro: tutte sono chiamate REST sincrone.

\Needspace{30\baselineskip}
\begin{figure}[H]
\centering
\begin{tikzpicture}[
    cmp/.style={rectangle, draw=#1!70!black, fill=#1!14, line width=0.8pt, font=\scriptsize, align=center,
                minimum width=2.4cm, minimum height=1.1cm, inner sep=2pt},
    gw/.style={cmp=swgreen},
    ui/.style={cmp=swblue},
    be/.style={cmp=sworange},
    chiam/.style={-{Straight Barb[length=2mm]}, dashed, line width=0.7pt, draw=black!70},
    api/.style={font=\tiny, inner sep=1.5pt, align=center},
    team/.style={rectangle, draw=sworange!70!black, fill=sworange!28, font=\tiny\bfseries, inner sep=1.5pt,
                 minimum width=0.75cm, anchor=south west}]
  \node[gw, minimum width=11.0cm, minimum height=0.9cm] (uig) at (5.8,0) {«component»\\ \textbf{UI Gateway}};
  \node[ui] (fe) at (1.5,-2.2)  {«component»\\ \textbf{Game Front End}};
  \node[ui] (st) at (5.8,-2.2)  {«component»\\ \textbf{Student Service}};
  \node[ui] (ad) at (10.1,-2.2) {«component»\\ \textbf{Admin Service}};
  \node[be] (re) at (1.5,-4.6)  {«component»\\ \textbf{Game Repository}};
  \node[be] (en) at (1.5,-7.0)  {«component»\\ \textbf{Game Engine}};
  \node[gw, minimum width=2.0cm] (apig) at (-3.5,-7.0) {«component»\\ \textbf{API Gateway}};
  \node[be] (ev) at (1.5,-9.4)  {«component»\\ \textbf{EvoSuite Runner}};
  \node[be] (tr) at (5.8,-9.4)  {«component»\\ \textbf{Student Test Runner}};
  \foreach \n/\t in {fe/T5, st/T2-3, ad/T1, re/T4, en/T6, ev/T8, tr/T7}
    \node[team] at (\n.north west) {\t};
  \foreach \n in {uig, fe, st, ad, re, en, apig, ev, tr} \compicon{\n};

  % l'UI Gateway serve le pagine dei tre servizi dell'interfaccia
  \foreach \n in {fe, st, ad}
    \draw[chiam] (uig.south -| \n) -- node[api, right]{HTTP GET} (\n.north);
  % chiamate fra i servizi dell'interfaccia
  \draw[chiam] ([yshift=0.25cm]fe.east) -- node[api, above]{HTTP GET} ([yshift=0.25cm]st.west);
  \draw[chiam] ([yshift=0.25cm]ad.west) -- node[api, above]{GET via\\ UI Gateway} ([yshift=0.25cm]st.east);
  \draw[chiam] ([yshift=-0.3cm]fe.east) -- ++(0.8,0) |- (9.6,-3.3)
    node[api, below, pos=0.75]{GET via UI Gateway} -- (9.6,-2.75);
  % verso il repository del gioco
  \draw[chiam] (fe.south) -- node[api, left]{GET /robots\\ GET /turns\\ (via UI Gateway)} (re.north);
  \draw[chiam] (10.6,-2.75) |- node[api, above, pos=0.75]{GET via UI Gateway, POST /robots} (re.east);
  \draw[chiam] (en.north) -- node[api, right]{GET /robot\_coverage\\ POST /turn\\ POST /close\_game} (re.south);
  % il motore di gioco
  \draw[chiam] ([yshift=0.25cm]en.east) -| node[api, above, pos=0.25]{GET /class\_ut\_code} (11.1,-2.75);
  \draw[chiam] ([yshift=-0.25cm]en.east) -| node[api, right, pos=0.75]{POST /compile\_and\_coverage} (tr.north);
  % l'API Gateway
  \draw[chiam] (fe.west) -| node[api, left, pos=0.75]{POST /compile\_and\_coverage\\ GET /class\_ut} (apig.north);
  \draw[chiam] (apig.east) -- node[api, below]{GET /class\_ut} (en.west);
  \draw[chiam] (apig.south) |- node[api, below right, pos=0.5]{/compute\_coverage\_metrics} (ev.west);
\end{tikzpicture}
\caption{Il diagramma dei componenti (ridisegnato e semplificato dalle slide). Le frecce tratteggiate sono chiamate REST,
         con il metodo HTTP e l'endpoint invocato; le etichette T1, \ldots, T8 indicano il team che ha sviluppato il
         servizio.}
\end{figure}

\subsection{Gli anti-pattern individuati}

Analizzando il sistema sono stati individuati molti \term{anti-pattern}: soluzioni ricorrenti che sembrano ragionevoli ma
producono più problemi di quanti ne risolvano. Il numero accanto al nome è quello dell'anti-pattern nel catalogo usato per
l'analisi; l'ultima colonna dice in quale servizio è stato trovato, oppure se riguarda \textbf{tutto il sistema}. Le slide
li raggruppano in quattro famiglie.

\Needspace{14\baselineskip}
\begin{center}
\renewcommand{\arraystretch}{1.3}
\begin{tabular}{@{}>{\bfseries\raggedright\arraybackslash}p{3.7cm} >{\raggedright\arraybackslash}p{9.6cm}
                >{\raggedright\arraybackslash}p{2.1cm}@{}}
  \toprule
  \multicolumn{3}{@{}l}{\textcolor{swteal}{\textbf{Scomposizione}}} \\
  \normalfont\textbf{Anti-pattern} & \textbf{Che cosa succede} & \textbf{Dove} \\
  \midrule
  \multicolumn{3}{@{}l}{\textit{Progettazione interna dei servizi}} \\
  No API-versioning (n.~6) & il sistema espone le sue API senza alcuna strategia di \textbf{versionamento}: ogni modifica
    rischia di rompere chi le usa & tutto il sistema \\
  Ambiguous service (n.~9) & le \textbf{responsabilità} di un servizio non sono definite con chiarezza: confusione e
    funzionalità sovrapposte & T23 \\
  \midrule
  \multicolumn{3}{@{}l}{\textit{Scomposizione fra servizi}} \\
  Chatty service (n.~17) & un servizio comunica \textbf{troppo} con gli altri, con molte chiamate a grana fine per
    raccogliere i dati che gli servono: cresce il tempo di risposta e peggiorano le prestazioni dell'intero sistema & T6 \\
  \bottomrule
\end{tabular}
\end{center}

Sono problemi di \textbf{confini}. Senza versioni delle API viene meno il vantaggio principale dei microservizi, il
rilascio indipendente: un servizio non può cambiare senza costringere gli altri ad adeguarsi. Un servizio dalle
responsabilità ambigue e uno che passa il tempo a interrogare gli altri sono il sintomo di bounded context tracciati male:
dati e operazioni che stanno insieme sono finiti in servizi diversi. Non a caso il servizio chiacchierone è il motore di
gioco (T6), che nel diagramma dei componenti chiama quasi tutti.

\Needspace{16\baselineskip}
\begin{center}
\renewcommand{\arraystretch}{1.3}
\begin{tabular}{@{}>{\bfseries\raggedright\arraybackslash}p{3.7cm} >{\raggedright\arraybackslash}p{9.6cm}
                >{\raggedright\arraybackslash}p{2.1cm}@{}}
  \toprule
  \multicolumn{3}{@{}l}{\textcolor{swteal}{\textbf{Interazione fra servizi}}} \\
  \normalfont\textbf{Anti-pattern} & \textbf{Che cosa succede} & \textbf{Dove} \\
  \midrule
  On-line only (n.~25) & il sistema si affida \textbf{solo} alla comunicazione sincrona fra microservizi ed evita
    l'elaborazione batch anche per i compiti lunghi o pesanti: colli di bottiglia e inefficienze & T1 \\
  Hardcoded endpoint (n.~28) & indirizzi, porte o endpoint dei microservizi sono \textbf{scritti} nel codice, nei file di
    configurazione o nelle variabili d'ambiente: tra l'altro, è difficile scalare & T1, T5, T6 \\
  No API gateway (n.~29) & manca un API gateway centralizzato: i client devono chiamare \textbf{direttamente} i servizi,
    e sicurezza e instradamento si complicano & tutto il sistema \\
  Wobbly service interactions (n.~30) & un servizio interagisce con altri servizi, o con i router dei messaggi, senza
    meccanismi per \textbf{tollerarne i guasti} & tutto il sistema \\
  Timeout (n.~31) & chi usa un servizio spesso non riesce a connettersi, probabilmente perché la sua
    \textbf{indisponibilità} non è gestita come si deve & tutto il sistema \\
  \bottomrule
\end{tabular}
\end{center}

Qui ritornano quasi tutti i temi del capitolo. \textit{On-line only} è esattamente il problema che le \textbf{code di
messaggi} risolvono: un compito lungo svolto in linea fa aspettare l'utente. Con gli endpoint scritti nel codice non si
può mettere un \textbf{load balancer} davanti a più istanze di un servizio, cioè non si può scalare orizzontalmente. Il
gateway, nei diagrammi, c'è; ma non è l'\textbf{unico} punto d'ingresso, visto che diverse chiamate lo scavalcano e
alcuni servizi sono raggiungibili direttamente dall'esterno. Infine \textit{Wobbly service interactions} e
\textit{Timeout} riguardano la \textbf{disponibilità}: in un sistema distribuito i guasti parziali sono la normalità, e un
servizio che non sa reagire all'assenza di un altro trascina con sé tutto il sistema.

\Needspace{22\baselineskip}
\begin{center}
\renewcommand{\arraystretch}{1.3}
\begin{tabular}{@{}>{\bfseries\raggedright\arraybackslash}p{3.7cm} >{\raggedright\arraybackslash}p{9.6cm}
                >{\raggedright\arraybackslash}p{2.1cm}@{}}
  \toprule
  \multicolumn{3}{@{}l}{\textcolor{swteal}{\textbf{Sicurezza}}} \\
  \normalfont\textbf{Anti-pattern} & \textbf{Che cosa succede} & \textbf{Dove} \\
  \midrule
  Unauthenticated traffic (n.~33) & le richieste alle API, sia dei client esterni sia fra microservizi interni, non sono
    \textbf{autenticate} come si deve & tutto il sistema \\
  Multiple user authentication (n.~34) & il sistema offre \textbf{più punti} di autenticazione degli utenti: aumentano la
    superficie d'attacco e i rischi & T1, T23 \\
  Publicly accessible microservices (n.~35) & i microservizi sono esposti \textbf{direttamente} ai client esterni: senza
    controlli adeguati possono rivelare informazioni sensibili, e il rischio di attacchi cresce & T1, T23, T4, T8 \\
  Unnecessary privileges to microservices (n.~36) & i microservizi hanno più accessi, permessi o funzionalità di quanti
    ne servano al loro compito & tutto il sistema \\
  Insufficient access control (n.~37) & manca un \textbf{controllo degli accessi} adeguato: dati sensibili e funzioni di
    business restano esposti e la riservatezza è compromessa & tutto il sistema \\
  Non-secured service-to-service communications (n.~39) & le interazioni fra servizi non sono \textbf{cifrate} né
    protette in altro modo: i dati si possono intercettare, alterare, spiare & tutto il sistema \\
  Hardcoded secrets (n.~42) & dati di configurazione sensibili, come password o chiavi delle API, sono scritti
    \textbf{direttamente} nel codice sorgente o negli script di deploy & T1, T23, T4 \\
  \bottomrule
\end{tabular}
\end{center}

Un monolite ha un solo ingresso da proteggere; un sistema a microservizi ne ha potenzialmente uno per servizio, più tutte
le comunicazioni interne, che viaggiano in rete. Molti di questi anti-pattern hanno la stessa cura: un \textbf{gateway}
che sia davvero l'unico punto d'ingresso, con un'\textbf{autenticazione centralizzata}, servizi raggiungibili solo dalla
rete interna, comunicazioni cifrate e segreti custoditi fuori dal codice.

\Needspace{24\baselineskip}
\begin{center}
\renewcommand{\arraystretch}{1.3}
\begin{tabular}{@{}>{\bfseries\raggedright\arraybackslash}p{3.7cm} >{\raggedright\arraybackslash}p{9.6cm}
                >{\raggedright\arraybackslash}p{2.1cm}@{}}
  \toprule
  \multicolumn{3}{@{}l}{\textcolor{swteal}{\textbf{Organizzazione dei team}}} \\
  \normalfont\textbf{Anti-pattern} & \textbf{Che cosa succede} & \textbf{Dove} \\
  \midrule
  Golden hammer (n.~44) & una tecnologia familiare viene applicata a \textbf{ogni} problema, adatta o no: gli
    sviluppatori riusano quella che conoscono bene anche dove non serve & tutto il sistema \\
  Too many standards (n.~46) & si usano troppi linguaggi, protocolli o framework diversi. I microservizi permettono la
    diversità tecnologica, ma un'eterogeneità \textbf{eccessiva} rende difficile la manutenzione, soprattutto quando gli
    sviluppatori cambiano, e complica integrazione ed evoluzione & T1, T23, T4, T8 \\
  Inadequate techniques support (n.~47) & il team usa strumenti, processi o metodi poco adatti a costruire e gestire
    microservizi, spesso pur avendo a disposizione alternative migliori & tutto il sistema \\
  No CI/CD (n.~53) & non si adottano l'\textbf{integrazione continua} e la consegna (o il rilascio) continui & tutto il
    sistema \\
  Manual configuration (n.~54) & host, servizi e istanze sono configurati \textbf{a mano}: non c'è una separazione netta
    fra il codice e la gestione della configurazione, e l'automazione è ostacolata & tutto il sistema \\
  Insufficient monitoring (n.~55) e Local logging (n.~58) & prestazioni e guasti dei microservizi non sono
    \textbf{monitorati}, quindi i problemi sono difficili da scoprire e da risolvere; e ogni servizio gestisce i propri log
    per conto suo, senza aggregarli in un punto unico & tutto il sistema \\
  Dismiss documentation (n.~56) & le API esposte non sono \textbf{documentate} a dovere: diventa difficile avere una
    visione d'insieme del sistema & tutto il sistema \\
  \bottomrule
\end{tabular}
\end{center}

\textit{Too many standards} è l'altra faccia di uno dei vantaggi dei microservizi, la libertà di scelta delle tecnologie,
e qui il ricambio degli sviluppatori è estremo: a ogni edizione del corso arrivano studenti nuovi. Senza CI/CD e con la
configurazione fatta a mano non si possono migliorare le metriche \textbf{DORA}, che erano la ragione stessa per adottare i
microservizi. Senza monitoraggio e senza log centralizzati manca l'\textbf{osservabilità} su cui si basa il fail-over.

\info{Che cosa insegna questo caso}{
  Quasi tutti gli anti-pattern riguardano \textbf{tutto il sistema}, non un singolo servizio: versioni delle API, gateway,
  autenticazione, tolleranza ai guasti, CI/CD, log centralizzati. Sono le decisioni che stanno \textbf{fra} i servizi, e
  che nessun team, sviluppando il proprio servizio per conto suo (qui addirittura in edizioni diverse del corso), aveva il
  compito di prendere. È la «complessità operativa» del prezzo dei microservizi, e insieme un'illustrazione della legge di
  Conway, a cui è dedicata la prossima sezione: team che non comunicano producono servizi che non si integrano.
}

% ============================================================
\section{Architettura e organizzazione dei team}

Che cosa c'entra l'organizzazione dei team con l'architettura? Moltissimo.

\dfn{Legge di Conway}{
  \textit{«Qualunque organizzazione che progetta un sistema produrrà un progetto la cui struttura è una copia della
  struttura di comunicazione dell'organizzazione.»}

  \hfill --- Melvin Conway, \textit{How do committees invent?}, Datamation, 1968
}

Il motivo è che l'accoppiamento del software è reso possibile e incoraggiato dalla \textbf{comunicazione fra persone}. Se si
può parlare con l'autore di un pezzo di codice, è più facile capire davvero quel codice; e quindi è più facile che il
proprio codice interagisca con esso, cioè che vi si \textbf{accoppi}. Due team che non si parlano produrranno, quasi
inevitabilmente, due moduli con un'interfaccia stretta fra loro.

\subsection{Come si organizza un team grande}

\begin{figure}[H]
\centering
\begin{tikzpicture}[
    team/.style={rectangle, rounded corners=3pt, draw=#1!70!black, fill=#1!12, line width=0.8pt, font=\scriptsize\bfseries,
                 minimum width=3.4cm, minimum height=0.65cm, align=center},
    liv/.style={rectangle, draw=#1!70!black, fill=#1!12, line width=0.8pt, font=\scriptsize,
                minimum width=3.6cm, minimum height=0.65cm}]
  % team per livello tecnologico
  \node[font=\small\bfseries] at (2.6,1.0) {Team per livello tecnologico};
  \node[team=q4] at (0.5,0.1) {\iuser\iuser\iuser\ specialisti UI/UX};
  \node[team=q7] at (0.5,-0.75) {\iuser\iuser\iuser\ middleware};
  \node[team=q5] at (0.5,-1.6) {\iuser\iuser\iuser\ DB admin};
  \node[liv=q4] (l1) at (5.0,0.1) {interfaccia};
  \node[liv=q7] (l2) at (5.0,-0.75) {logica};
  \node[liv=q5] (l3) at (5.0,-1.6) {dati};
  \foreach \y in {0.1,-0.75,-1.6} \draw[fr] (2.3,\y) -- (3.1,\y);
  \node[note, align=center, text=swred] at (2.6,-2.45) {nuova funzionalità: tre team da coordinare,\\ logica sparsa ovunque};
  % team trasversali
  \begin{scope}[shift={(8.6,0)}]
    \node[font=\small\bfseries] at (2.6,1.0) {Team trasversali (cross-functional)};
    \foreach \t/\y/\n in {A/0.1/Ordini, B/-0.75/Catalogo, C/-1.6/Pagamenti} {
      \node[team=swgreen, minimum width=2.6cm] at (0.6,\y) {\iuser\iuser\iuser\ Team \t};
      \node[servizio, minimum width=2.2cm, minimum height=0.6cm] at (4.4,\y) {\n};
      \draw[fr] (2.0,\y) -- (3.2,\y);
    }
    \node[note, align=center, text=swgreen!70!black] at (2.6,-2.45) {ogni team ha tutte le competenze\\ e possiede un servizio};
  \end{scope}
\end{tikzpicture}
\caption{La legge di Conway all'opera: team divisi per tecnologia producono un'architettura a livelli; team trasversali
         producono servizi organizzati attorno alle capacità di business.}
\end{figure}

Spesso il management organizza i team in base al \textbf{livello tecnologico}: specialisti di interfaccia, di middleware, di
database. Team funzionali divisi per livello portano ad \textbf{architetture a livelli}. E se serve aggiungere una nuova
funzionalità? Bisogna coordinare tutti i team, e la logica finisce un po' ovunque.

I microservizi favoriscono invece team \textbf{piccoli e indipendenti}, e \textbf{trasversali} (\textit{cross-functional}):
ogni team include l'intera gamma di competenze necessarie (interfaccia, logica, dati, test, rilascio) e possiede uno o più
servizi.

\subsection{Che fare con la legge di Conway}

\begin{itemize}
  \item \textbf{Ignorarla}: non tenerne conto, perché non se ne è mai sentito parlare o perché si pensa che non valga nel
        proprio caso. \textit{Spoiler}: vale.
  \item \textbf{Accettarla}: riconoscerne l'impatto, e fare in modo che l'architettura non si scontri con i canali di
        comunicazione del team.
  \item \textbf{La manovra di Conway inversa} (\textit{Inverse Conway Maneuver}): cambiare la struttura di comunicazione del
        team per \textbf{favorire l'architettura desiderata}. Se vogliamo tre servizi indipendenti, organizziamo tre team
        indipendenti.
\end{itemize}

\info{Il messaggio da portare a casa}{
  La scomposizione di un sistema e la scomposizione dell'organizzazione che lo sviluppa \textbf{non si possono fare
  indipendentemente}: né all'inizio del progetto, né durante tutto il suo ciclo di vita. Far evolvere l'architettura e
  riorganizzare le persone devono andare di pari passo.
}

% ============================================================
% il riepilogo e' lungo: deve stare tutto in una pagina
\Needspace{30\baselineskip}
\gbox{In sintesi}{azure-gradient-6!80!black}{
\begin{itemize}
  \item L'architettura è l'insieme delle \textbf{decisioni principali}: quelle difficili da cambiare in seguito. Una buona
        qualità interna costa all'inizio, ma ripaga in poche settimane.
  \item \textbf{Scalabilità}: le prestazioni crescono in proporzione alle risorse. Si scala orizzontalmente,
        verticalmente o diagonalmente; per scalare orizzontalmente serve un \textbf{load balancer}, che non deve diventare
        un collo di bottiglia o un single point of failure.
  \item \textbf{Disponibilità}: si misura in «nove» e si ottiene con il \textbf{fail-over} attivo-passivo o attivo-attivo,
        con heartbeat e osservabilità per accorgersi dei guasti.
  \item \textbf{Dati}: le proprietà \textbf{ACID} rendono difficile scalare un RDBMS; le tecniche sono replica
        (master-slave, master-master), federazione, sharding, denormalizzazione. I database \textbf{NoSQL} (BASE) scalano
        meglio rinunciando alla consistenza immediata: \textbf{consistenza finale}.
  \item Le \textbf{code di messaggi} rendono asincrone le elaborazioni lente e riducono la latenza; un broker serve anche
        a consegnare lo stesso evento a più consumatori, per esempio a tutti i dispositivi di un utente.
  \item I \textbf{microservizi} danno confini forti, rilasci indipendenti e libertà tecnologica, al prezzo di complessità,
        consistenza finale e costi operativi: vanno adottati solo per sistemi complessi, progettandone i confini con il
        DDD (bounded context, aggregati).
  \item Il \textbf{caso reale} di un gioco a microservizi mostra gli \textbf{anti-pattern} tipici (servizi chiacchieroni,
        endpoint scritti nel codice, nessun gateway unico, autenticazione sparsa, niente CI/CD né log centralizzati):
        quasi tutti riguardano ciò che sta \textbf{fra} i servizi.
  \item \textbf{Legge di Conway}: la struttura del sistema copia quella della comunicazione dell'organizzazione. Architettura
        e organizzazione dei team vanno progettate insieme.
\end{itemize}
}
